\chapter{Tinjauan Pustaka}

\section{Tinjauan Pustaka}

Penelitian-penelitian yang sangat berkaitan dan menginspirasi penelitian. Dasar-dasar penelitian sebelumnya yang menjadi tinjauan pustaka pada penelitian ini dirangkum dalam Tabel \ref{tab:tinjauan-pustaka}.

\begin{table}[h]
\centering
\caption{Tinjauan pustaka}
\label{tab:tinjauan-pustaka}
\begin{tabular}{|p{3cm}|p{3cm}|p{2.5cm}|p{2.5cm}|p{2.5cm}|}
\hline
\textbf{Author} & \textbf{Judul Penelitian} & \textbf{Metode} & \textbf{Hasil} & \textbf{Kelemahan} \\
\hline
Contoh Author 1 & Judul penelitian pertama & Metode yang digunakan & Hasil yang diperoleh & Kelemahan penelitian \\
\hline
Contoh Author 2 & Judul penelitian kedua & Metode yang digunakan & Hasil yang diperoleh & Kelemahan penelitian \\
\hline
Contoh Author 3 & Judul penelitian ketiga & Metode yang digunakan & Hasil yang diperoleh & Kelemahan penelitian \\
\hline
\end{tabular}
\end{table}

\textcolor{red}{Catatan: Ganti contoh di atas dengan penelitian-penelitian yang relevan dengan topik Anda. Tambahkan baris sesuai kebutuhan. Pastikan setiap penelitian yang disebutkan di tabel juga dicantumkan di daftar pustaka dan disitasi dengan benar.}

Setelah membahas berbagai penelitian terdahulu melalui tabel di atas, dapat disimpulkan bahwa masih terdapat peluang pengembangan atau gap penelitian yang akan diisi oleh penelitian ini.
